
\section{Conclusions and Future Work}
\label{sec:conclusion}

\noindent
Our experiments show that sketchiness can be used as an independent visual variable to map ordinal data in 2.5D treemaps without having an impact on reading and interpreting other visual variables such as color and height. 
The general weakness of sketchy rendering, 
the lowered data/ink ratio, 
is less relevant due to the uniform and regular block shape used by 2.5D treemaps---%
there is no noticeable \enquote{chart junk} \cite{Tufte1986} introduced.
Additionally, we could measure a high accuracy for the selective pre-attentive processing for 2 of 3 styles (Marker, Pencil). 
The use of a conjunction of sketchiness and color is poorly processed pre-attentively, as expected. 
In addition, the pure use of outlines (as in the Biro style) leads to poor pre-attentive processing 
and as the size of elements increases, the ability to pre-attentively process the items decreases. 
With the length of the design space of 3 to 4 distinct values as described by Boukhelifa et al., we could also show that users are able to bring depicted elements into an order 
based on their sketchiness value for all styles. 
In summary, the expressiveness of 2.5D treemaps can be extended by sketchiness as visual variable in its own right if applied to appropriate data such as ordinal data with small range. 

In future work, we would like to explore how time-varying uncertainty data can be mapped and animated in 2.5D treemaps, \eg to explore changes over time. 
We also want to explore how sketchy outlines can be integrated into a global illumination concept for 2.5D treemaps (e.g., shadow casting) and how sketchy line patterns for general ordinal data not related to uncertainty could look like.

